\begin{table}[ht]
\centering\footnotesize
\begin{tabular}{@{}rrllc@{}}
\toprule
$p$ & $s_0$ & $(a_1,a_2,a_3,a_4)$ & $P_{\mathfrak p} \bmod 2$ & order\\
\midrule
$7$ & $5$ & $(-1, -2, 2, 38)$ & $(x+1)^{2}(x^{3}+x+1)(x^{3}+x^{2}+1)$ & $7$\\
$19$ & $3$ & $(-5, 38, -152, 650)$ & $(x+1)^{2}(x^{3}+x+1)(x^{3}+x^{2}+1)$ & $7$\\
$31$$^\dagger$ & $12$ & $(11, 88, 682, 4438)$ & $(x+1)^{2}(x^{3}+x+1)(x^{3}+x^{2}+1)$ & $7$\\
$41$ & $8$ & $(-2, 56, -128, 2402)$ & $(x+1)^{8}$ & $1$\\
$43$ & $15$ & $(-1, -10, -226, 1278)$ & $(x+1)^{2}(x^{3}+x+1)(x^{3}+x^{2}+1)$ & $7$\\
\bottomrule
\end{tabular}
\caption{Frobenius polynomials of the covers defined over $L=\QQ(s)$ at the primes $\mathfrak p=(p,s-s_0)$ of degree one, with the notation of Table~\ref{tab:frobenius}. In every case $P_{\mathfrak p}\equiv QQ^{\rev}\pmod 2$ and the congruences of Proposition~\ref{prop:23} hold. At the prime marked with a dagger the model has $57$ nodes, a tacnode and an ordinary triple point.}
\label{tab:lcovers}
\end{table}
